% answers_short.tex —— 2002 年·北京卷·文科 · 简短答案速查
% 由 tools/make_short_answers.py 自 answers.tex 的 \gksolhead/\gkanswer 提取，
% 并与 2002/2002_数学真题_北京文_答案要点.md 逐题比对后生成，勿手改（改源 MD 后重跑）。

\begin{gknotice}
\textbf{使用说明}：本节只给最终答案便于快速核对，完整推导见解析版。答案为编写者独立推导，\textbf{非官方标准答案}。
\end{gknotice}

\gkpart{选择题}
\begin{center}\begin{tabular}{ccccccccccccc}
\toprule
题号 & 1 & 2 & 3 & 4 & 5 & 6 & 7 & 8 & 9 & 10 & 11 & 12 \\
\midrule
答案 & C & D & B & A & C & B & D & A & B & D & C & A \\
\bottomrule
\end{tabular}\end{center}

\gkpart{填空题}
\begin{itemize}[leftmargin=2.6em]
  \item[\textbf{13.}] $\cos\frac{6\pi}{5}<\sin\frac{2\pi}{5}<\tan\frac{7\pi}{5}$。
  \item[\textbf{14.}] $4$。
  \item[\textbf{15.}] ①②③④⑤。
  \item[\textbf{16.}] $2$。
\end{itemize}

\gkpart{解答题}
\begin{itemize}[leftmargin=2.6em]
  \item[\textbf{17.}] $\left[\frac12,5\right)$。
  \item[\textbf{18.}] (1) $\tan\theta=\dfrac{2h}{b-d}$；(2) $V_{\mathrm{\text{估}}}<V$
  \item[\textbf{19.}] (1) $x_n\geqslant\sqrt a\quad(n\geqslant2)$；(2) $x_n\geqslant x_{n+1}\quad(n\geqslant2)$
  \item[\textbf{20.}] (1) 至少2个单位时间。第一轮：1读2、2读1、3读4、4读3；第二轮：1读3、2读4、3读1、4读2；第二轮后四台机器都得到总和；(2) 7个单位时间
  \item[\textbf{21.}] (1) $G(\frac{1+b}{3},\frac c3)$，$F(\frac12,\frac{b^2+c^2-b}{2c})$，$H(b,\frac{b(1-b)}c)$；三点共线；(2) $3\left(b-\frac12\right)^2+c^2=\frac34$，去掉$(0,0)$、$(1,0)$及$(\frac12,\pm\frac{\sqrt3}{2})$
  \item[\textbf{22.}] (1) $f(0)=f(1)=0$；(2) $f$为奇函数；(3) $u_n=n2^n$，故$u_{n+1}-u_n=(n+2)2^n>0$
\end{itemize}
